\section{Ninja's Training}
\label{sec:dp-ninja-training}

\problemstatement{A ninja has $n$ days of training. On each day, they
must perform exactly one of $3$ activities (indexed $0,1,2$), earning
$\textit{points}[\textit{day}][\textit{activity}]$ points -- but they may
\textbf{not} perform the \textbf{same} activity on two consecutive days.
Maximize the total points earned over all $n$ days.}

\begin{patternbox}
\emph{``Choose one of several options each day, but not the same option
as the previous day''} is the keyword pattern from
Section~\ref{sec:dp-grids-intro} for a DP state with a \textbf{non-spatial
second dimension}: the state needs both \emph{which day} we are on and
\emph{which activity was excluded} (because it was used the day before),
even though there is no grid or spatial structure here at all -- the
``two dimensions'' are day and constraint, not row and column.
\end{patternbox}

\subsection*{Deriving the recurrence}

Let $f(\textit{day}, \textit{last})$ be the maximum points earnable from
day $0$ through day $\textit{day}$, given that activity $\textit{last}$
is forbidden on day $\textit{day}$ specifically (because it was used on
day $\textit{day}+1$, which we are building backward from -- or, viewed
forward, because it was the most recently committed activity choice we
must avoid repeating). For each candidate activity $\textit{task} \ne
\textit{last}$:
\[
  f(\textit{day}, \textit{last}) = \max_{\textit{task} \ne \textit{last}}
  \Big(\textit{points}[\textit{day}][\textit{task}] +
  f(\textit{day}-1, \textit{task})\Big).
\]
Base case: $f(0, \textit{last}) = \max_{\textit{task} \ne \textit{last}}
\textit{points}[0][\textit{task}]$ (no earlier day to recurse into). The
final answer is $f(n-1, 3)$, using a sentinel value $3$ (outside the
valid activity range $\{0,1,2\}$) to mean ``no activity is forbidden on
the last day,'' since nothing comes after it.

\subsection*{Approach 1: Recursion}

\begin{lstlisting}
#include <bits/stdc++.h>
using namespace std;

int ninjaTrainingRecursive(int day, int last, vector<vector<int>>& points) {
    if (day == 0) {
        int best = 0;
        for (int task = 0; task < 3; task++) {
            if (task != last) best = max(best, points[0][task]);
        }
        return best;
    }

    int best = 0;
    for (int task = 0; task < 3; task++) {
        if (task != last) {
            int candidate = points[day][task] +
                ninjaTrainingRecursive(day - 1, task, points);
            best = max(best, candidate);
        }
    }
    return best;
}

int main() {
    vector<vector<int>> points = {{10,40,70},{20,50,80},{30,60,90}};
    int n = points.size();
    cout << ninjaTrainingRecursive(n - 1, 3, points) << endl; // 210
    return 0;
}
\end{lstlisting}

\begin{complexitybox}[Approach 1 Complexity]
\textbf{Time.} $O(3^n)$ -- up to $3$ branches per call, depth $n$.
\textbf{Space.} $O(n)$ recursion stack.
\end{complexitybox}

\subsection*{Approach 2: Memoization}

\begin{lstlisting}
#include <bits/stdc++.h>
using namespace std;

int ninjaTrainingMemo(int day, int last, vector<vector<int>>& points,
                       vector<vector<int>>& memo) {
    if (day == 0) {
        int best = 0;
        for (int task = 0; task < 3; task++) {
            if (task != last) best = max(best, points[0][task]);
        }
        return best;
    }
    if (memo[day][last] != -1) return memo[day][last];

    int best = 0;
    for (int task = 0; task < 3; task++) {
        if (task != last) {
            int candidate = points[day][task] +
                ninjaTrainingMemo(day - 1, task, points, memo);
            best = max(best, candidate);
        }
    }
    return memo[day][last] = best;
}

int main() {
    vector<vector<int>> points = {{10,40,70},{20,50,80},{30,60,90}};
    int n = points.size();
    vector<vector<int>> memo(n, vector<int>(4, -1));
    cout << ninjaTrainingMemo(n - 1, 3, points, memo) << endl; // 210
    return 0;
}
\end{lstlisting}

\begin{complexitybox}[Approach 2 Complexity]
\textbf{Time.} $O(n \cdot 4 \cdot 3) = O(n)$ -- $n$ days, $4$ possible
\textit{last} values, $3$ choices of \textit{task} per state.
\textbf{Space.} $O(n)$ memo (the $4$ factor is a constant) $+$ $O(n)$
recursion stack.
\end{complexitybox}

\subsection*{Approach 3: Tabulation}

\begin{lstlisting}
#include <bits/stdc++.h>
using namespace std;

int ninjaTrainingTabulation(vector<vector<int>>& points) {
    int n = points.size();
    vector<vector<int>> dp(n, vector<int>(4, 0));

    for (int last = 0; last < 4; last++) {
        int best = 0;
        for (int task = 0; task < 3; task++) {
            if (task != last) best = max(best, points[0][task]);
        }
        dp[0][last] = best;
    }

    for (int day = 1; day < n; day++) {
        for (int last = 0; last < 4; last++) {
            int best = 0;
            for (int task = 0; task < 3; task++) {
                if (task != last) {
                    int candidate = points[day][task] + dp[day-1][task];
                    best = max(best, candidate);
                }
            }
            dp[day][last] = best;
        }
    }
    return dp[n - 1][3];
}

int main() {
    vector<vector<int>> points = {{10,40,70},{20,50,80},{30,60,90}};
    cout << ninjaTrainingTabulation(points) << endl; // 210
    return 0;
}
\end{lstlisting}

\subsection*{Dry run of the tabulation table}

\dryrun{Take $\textit{points} = [[10,40,70],[20,50,80],[30,60,90]]$.}

\begin{center}
\begin{tabular}{c|c|c|c|c}
$\textit{day}$ & $\textit{last}=0$ & $\textit{last}=1$ &
$\textit{last}=2$ & $\textit{last}=3$ \\ \hline
$0$ & 70 & 70 & 40 & 70 \\ \hline
$1$ & 120 & 120 & 120 & 120 \\ \hline
$2$ & 210 & 210 & 180 & 210
\end{tabular}
\end{center}

Reading $\textit{dp}[0][2]=40$: with activity $2$ forbidden on day $0$,
the best of $\{\textit{points}[0][0],\textit{points}[0][1]\} =
\{10,40\}$ is $40$. Reading $\textit{dp}[1][3]$ (no restriction on day
$1$): $\max(20+\textit{dp}[0][0],\ 50+\textit{dp}[0][1],\
80+\textit{dp}[0][2]) = \max(20{+}70,\ 50{+}70,\ 80{+}40) =
\max(90,120,120) = 120$. Reading the final answer,
$\textit{dp}[2][3]$: $\max(30+\textit{dp}[1][0],\ 60+\textit{dp}[1][1],\
90+\textit{dp}[1][2]) = \max(30{+}120,\ 60{+}120,\ 90{+}120) =
\max(150,180,210) = 210$.

\begin{complexitybox}[Approach 3 Complexity]
\textbf{Time.} $O(n)$ (the $4 \times 3$ factor is a constant).
\textbf{Space.} $O(n)$ table, no recursion stack.
\end{complexitybox}

\subsection*{Approach 4: Space Optimization}

\begin{lstlisting}
#include <bits/stdc++.h>
using namespace std;

int ninjaTrainingSpaceOptimized(vector<vector<int>>& points) {
    int n = points.size();
    vector<int> prevDay(4, 0);

    for (int last = 0; last < 4; last++) {
        int best = 0;
        for (int task = 0; task < 3; task++) {
            if (task != last) best = max(best, points[0][task]);
        }
        prevDay[last] = best;
    }

    for (int day = 1; day < n; day++) {
        vector<int> currentDay(4, 0);
        for (int last = 0; last < 4; last++) {
            int best = 0;
            for (int task = 0; task < 3; task++) {
                if (task != last) {
                    best = max(best, points[day][task] + prevDay[task]);
                }
            }
            currentDay[last] = best;
        }
        prevDay = currentDay;
    }
    return prevDay[3];
}

int main() {
    vector<vector<int>> points = {{10,40,70},{20,50,80},{30,60,90}};
    cout << ninjaTrainingSpaceOptimized(points) << endl; // 210
    return 0;
}
\end{lstlisting}

\begin{complexitybox}[Approach 4 Complexity]
\textbf{Time.} $O(n)$.
\textbf{Space.} $O(1)$ (a fixed-size array of $4$ entries, independent
of $n$).
\end{complexitybox}

\begin{takeawaybox}
A DP state's second dimension does not need to be spatial -- here, it is
``which choice is currently forbidden,'' a small, fixed-size constraint
dimension exactly analogous to the visited-tracking state used in
backtracking, but folded directly into the DP table's indexing instead
of an external visited structure. Whenever a problem's choice at step
$i$ is constrained only by the choice at step $i-1$ (not by anything
further back), this ``day $\times$ last choice'' state shape applies
directly.
\end{takeawaybox}
